\clearpage
\begin{explanationbox}
\textbf{一、为什么会想到“周期”？从一串数开始看。}
\begin{center}\begin{tikzpicture}[font=\small,>=stealth,line join=round]
\draw[gray!40,rounded corners=3pt] (0,0) rectangle (14.2,8.3);\node at (7.1,7.700000000000001) {例1：每走2，高度在3与负3之间交替};
\begin{scope}[shift={(1.5444444444444445,3.828947368421052)},x=0.5555555555555555cm,y=0.4842105263157894cm]\draw[->,gray] (-0.8,0)--(10,0) node[right]{$x$};\draw[->,gray] (0,-4.5)--(0,5) node[above]{$y$};\draw[blue,thick,solid,domain=-0.4:9.5,samples=180] plot(\x,{3*sin(deg(pi*\x/2))});\fill[orange!85!black] (1,3) circle(2.2pt);\fill[orange!85!black] (3,-3) circle(2.2pt);\fill[orange!85!black] (5,3) circle(2.2pt);\fill[orange!85!black] (7,-3) circle(2.2pt);\fill[orange!85!black] (9,3) circle(2.2pt);\foreach \v in {1,3,5,7,9}{\node[below=5pt] at (\v,0){$\v$};}\end{scope}\node[anchor=west,align=left,text width=5.0cm] at (8.4,4.0) {$f=3\sin\frac{\pi x}2$\\[9pt]$1,3,5,7,9$\\[9pt]对应 $3,-3,3,-3,3$\\[9pt]隔4回原值\\[9pt]$f(9)=f(1)=3$};
\node[text width=12.6cm,align=center] at (7.1,.65) {示例满足初值3；不能把每次相反误认为每次相等。};
\end{tikzpicture}\end{center}



先不去想曲线，就把函数看成“输入一个数、输出一个数”的机器。假设 $a=2$，某一点的函数值是 $3$，题目告诉我们：输入每增加 $2$，输出就变成上一次的相反数。
\[
\begin{array}{c|ccccc}
\text{输入}&1&3&5&7&9\\\hline
\text{输出}&3&-3&3&-3&3
\end{array}
\]
图中输入每增加2，高度就在3与负3之间交替。\textbf{“变号”不是“回到原值”；变号两次才是。}


把 $v$ 换成任意的 $f(x)$，仍然成立。这就解释了为什么\textbf{不需要知道具体的函数式子}，照样能找到一个周期。

\medskip
\end{explanationbox}\clearpage\begin{explanationbox}\textbf{二、图像上到底发生了什么？}
\begin{center}\begin{tikzpicture}[font=\small,>=stealth,line join=round]
\draw[gray!40,rounded corners=3pt] (0,0) rectangle (14.2,8.3);\node at (7.1,7.700000000000001) {隔2变为相反高度，隔4才回到原高度};
\begin{scope}[shift={(1.7153846153846155,3.8842857142857143)},x=0.7692307692307693cm,y=1.3142857142857143cm]\draw[->,gray] (-0.8,0)--(7,0) node[right]{$x$};\draw[->,gray] (0,-1.7)--(0,1.8) node[above]{$y$};\draw[blue,thick,solid,domain=-0.4:6.6,samples=180] plot(\x,{sin(deg(pi*\x/2))});\fill[orange!85!black] (1,1) circle(2.2pt);\fill[orange!85!black] (3,-1) circle(2.2pt);\fill[orange!85!black] (5,1) circle(2.2pt);\foreach \v/\h in {1/1,3/-1,5/1}{\draw[gray,dashed] (\v,0)--(\v,\h);\node[below=5pt] at (\v,0){$\v$};}\end{scope}\node[anchor=west,align=left,text width=5.0cm] at (8.4,4.0) {$f=\sin\frac{\pi x}2,\ a=2$\\[9pt]$f(1)=1$\\[9pt]$f(3)=-1$\\[9pt]$f(5)=1$\\[9pt]正周期可取4};
\node[text width=12.6cm,align=center] at (7.1,.65) {蓝线是满足条件的一条示例曲线；证明只用任意输入的两次变号。};
\end{tikzpicture}\end{center}



如果图像上有一点 $(x,y)$，那么一定还存在下面两点：
\[
(x+a,-y),\qquad (x+2a,y).
\]
也就是说：横坐标向右移动 $a$（当 $a>0$），纵坐标就上下颠倒；再向右移动 $a$，纵坐标回到原来的高度。它描述的是“\textbf{水平移动后，高低互换}”，\textbf{并不等于}“绕某个点转半圈”。

\textbf{看得见的例子：}取 $f(x)=\sin\!\frac{\pi x}{2}$。它满足 $f(x+2)=-f(x)$，所以 $4$ 是一个周期。

图中 $a=2$，同一高度的点隔 $4$ 重现；隔 $2$ 则高低相反。注意 $f(1)=1\ne0$，所以 $(1,0)$ 不是该图像的对称中心。

\medskip
\end{explanationbox}\clearpage\begin{explanationbox}\textbf{三、和已经学过的奇函数、偶函数联系起来（很值得记）。}
\begin{center}\begin{tikzpicture}[font=\small,>=stealth,line join=round]
\draw[gray!40,rounded corners=3pt] (0,0) rectangle (14.2,8.5);\node at (7.1,7.9) {同样加π变号，额外的奇偶性决定不同对称};
\node at (2.4,6.9) {$f=\sin x$：奇函数};\begin{scope}[shift={(1.2666666666666666,4.341935483870968)},x=0.7083333333333334cm,y=1.1612903225806452cm]\draw[->,gray] (-0.8,0)--(4,0) node[right]{$x$};\draw[->,gray] (0,-1.5)--(0,1.6) node[above]{$y$};\draw[blue,thick,solid,domain=-0.5:3.8,samples=160] plot(\x,{sin(deg(\x))});\fill[orange!85!black] (0.5235987755982988,0.5) circle(2pt);\fill[orange!85!black] (2.6179938779914944,0.5) circle(2pt);\draw[orange!85!black,dashed] (1.570796,-1.2)--(1.570796,1.3);\end{scope}\node[align=center,text width=4cm] at (2.4,1.7) {关于直线 $x=\frac\pi2$ 对称};\node at (7.050000000000001,6.9) {$g=\cos x$：偶函数};\begin{scope}[shift={(5.916666666666667,4.341935483870968)},x=0.7083333333333334cm,y=1.1612903225806452cm]\draw[->,gray] (-0.8,0)--(4,0) node[right]{$x$};\draw[->,gray] (0,-1.5)--(0,1.6) node[above]{$y$};\draw[violet,thick,solid,domain=-0.5:3.8,samples=160] plot(\x,{cos(deg(\x))});\fill[orange!85!black] (0.5235987755982988,0.8660254037844386) circle(2pt);\fill[orange!85!black] (2.6179938779914944,-0.8660254037844386) circle(2pt);\fill[orange!85!black] (1.570796,0) circle(2.2pt);\draw[orange!85!black,dashed] (.523599,.866025)--(2.617994,-.866025);\end{scope}\node[align=center,text width=4cm] at (7.050000000000001,1.7) {关于点 $(\frac\pi2,0)$\\中心对称};\node at (11.7,6.9) {$f=\sin x$ 的中心检验};\begin{scope}[shift={(10.566666666666666,4.341935483870968)},x=0.7083333333333334cm,y=1.1612903225806452cm]\draw[->,gray] (-0.8,0)--(4,0) node[right]{$x$};\draw[->,gray] (0,-1.5)--(0,1.6) node[above]{$y$};\draw[blue,thick,solid,domain=-0.5:3.8,samples=160] plot(\x,{sin(deg(\x))});\fill[orange!85!black] (1.5707963267948966,1) circle(2pt);\end{scope}\node[align=center,text width=4cm] at (11.7,1.7) {$f(\frac\pi2)=1\ne0$\\$(\frac\pi2,0)$ 不是中心};
\node[text width=12.6cm,align=center] at (7.1,.65) {三图同尺度；只有变号条件时，不能直接指定对称轴或对称中心。};
\end{tikzpicture}\end{center}



设 $c=\frac a2$。已知条件还可以写成
\[
f(c+t)=-f(t-c).
\]
接下来\textbf{必须另加条件}，才有对称结论：
\begin{itemize}[leftmargin=*,itemsep=0.5em]
\item \textbf{若还知道 $f$ 是偶函数：} $f(t-c)=f(c-t)$，于是 $f(c+t)=-f(c-t)$。两侧函数值相反，图像关于 $\bigl(\frac a2,0\bigr)$ \textbf{中心对称}。例如 $\cos x$、$a=\pi$。
\item \textbf{若还知道 $f$ 是奇函数：} $f(t-c)=-f(c-t)$，于是 $f(c+t)=f(c-t)$。两侧函数值相等，图像关于\textbf{直线 $x=\frac a2$ 对称}。例如 $\sin x$、$a=\pi$。
\end{itemize}
\textbf{惊喜在这里：}同样的“加 $a$ 变号”，配上\textbf{偶函数}得到一个对称中心，配上\textbf{奇函数}却得到一条对称轴。两种情形不能混淆！

\medskip
\end{explanationbox}\clearpage\begin{explanationbox}\textbf{四、还能从负号想到什么？}
\begin{center}\begin{tikzpicture}[font=\small,>=stealth,line join=round]
\draw[gray!40,rounded corners=3pt] (0,0) rectangle (14.2,8.3);\node at (7.1,7.700000000000001) {连续曲线：从正高度到负高度，途中必碰横轴};
\begin{scope}[shift={(1.7666666666666666,3.8842857142857143)},x=2.2222222222222223cm,y=1.3142857142857143cm]\draw[->,gray] (-0.3,0)--(2.4,0) node[right]{$x$};\draw[->,gray] (0,-1.7)--(0,1.8) node[above]{$y$};\draw[blue,thick,solid,domain=-0.1:2.1,samples=180] plot(\x,{cos(deg(pi*\x))});\fill[orange!85!black] (0,1) circle(2.2pt);\fill[orange!85!black] (1,-1) circle(2.2pt);\fill[orange!85!black] (0.5,0) circle(2.2pt);\end{scope}\node[anchor=west,align=left,text width=5.0cm] at (8.4,4.0) {$f=\cos\pi x,\ a=1$\\[9pt]$f(0)=1,\ f(1)=-1$\\[9pt]在 $[0,1]$ 连续\\[9pt]$f(\frac12)=0$\\[9pt]至少有一个零点};
\node[text width=12.6cm,align=center] at (7.1,.65) {连续是保证中途取到0的关键条件；端点本来为0时也算碰到横轴。};
\end{tikzpicture}\end{center}


\begin{itemize}[leftmargin=*,itemsep=0.4em]
\item \textbf{成对抵消：}$f(x)+f(x+a)=0$。因此 $f(x)+f(x+a)+f(x+2a)+f(x+3a)=0$，成千上万项也能两项一组计算。
\item \textbf{连续图像必过横轴：}若 $a>0$ 且 $f$ 在 $[x,x+a]$ 上连续，两端的函数值互为相反数，曲线必然在这段内与 $x$ 轴相交（包括端点）。\textbf{不连续时不能这么说。}
\item \textbf{负数 $a$ 也不怕：}假设 $a=-2$，由 $f(x-2)=-f(x)$，把 $x$ 改成 $x+2$，得到 $f(x+2)=-f(x)$，所以正周期仍可取 $4=2|a|$。
\end{itemize}
\end{explanationbox}
